\documentclass[conference,a4paper]{IEEEtran}

\begin{document}

\title{GAN-Based Heatwave-Oriented Temperature Imputation and Forecasting over Chennai}

\author{\IEEEauthorblockN{Author Name}
\IEEEauthorblockA{\textit{Department Name} \\
\textit{Institution Name}\\
City, Country \\
email address}
}

\maketitle

\begin{abstract}
Heatwaves pose a growing threat to coastal South Indian cities, and their analysis depends on long, continuous temperature records, which are often incomplete. We propose a generative adversarial network (GAN) based framework for heatwave-oriented temperature imputation and forecasting over the Chennai region. Long-term daily climate data for the period 1950 to 2025 were extracted from the ERA5-Land reanalysis through API-based retrieval from Google Earth Engine (GEE). Following exploratory data analysis and heatwave characterization, three models were developed and evaluated: a conventional GAN, a Shared GAN in which a single generator performs both imputation and forecasting, and a Shared Wasserstein GAN with Gradient Penalty (Shared WGAN-GP). The models were trained and tested on a chronological split with simulated random and multi-day gaps, and their performance was comparatively assessed for missing-data imputation and short-range temperature forecasting, with particular emphasis on heatwave events. Among the evaluated architectures, the Shared WGAN-GP demonstrated comparatively superior performance while requiring fewer parameters and less training time than separately trained task-specific networks, indicating its potential for robust climate-data imputation and heatwave forecasting. Future work will advance these imputations and forecasts towards El Ni\~no-aware modelling by incorporating El Ni\~no--Southern Oscillation indices, in order to capture the influence of large-scale climate variability on regional heatwaves.
\end{abstract}

\begin{IEEEkeywords}
El Ni\~no, generative adversarial networks, heatwaves, missing-data imputation, temperature forecasting, Wasserstein GAN
\end{IEEEkeywords}

\end{document}
